\chapter{Reproducibility and Data Audit}
\label{app:reproducibility}


% =================================================================
% A.1 Software and Hardware Environments
% =================================================================

\section{Software and Hardware Environments}
\label{app:environments}

All models trained for this thesis used an AutoDL-hosted NVIDIA
GeForce RTX 4090 instance. Local development (Windows with WSL2) was
CPU-only and used only for data-loading and pipeline debugging, never
for a reported training run. The published BiometryNet reference
numbers are not trained here and are excluded from this section.

\textbf{QCLand.} The core environment is pinned in
\path{requirements.txt}: \path{torch==2.7.0},
\path{torchvision==0.22.0}, \path{lightning==2.5.1.post0},
\path{timm==1.0.15}, \path{transformers==4.56.1},
\path{scipy==1.15.2}, \path{torchmetrics==1.7.1}, and
\path{jsonargparse==4.38}. This file records the declared
dependency specification rather than a verified snapshot of every
historical training run's exact resolved environment.

The following \texttt{timm} log lines confirm that constructing
\path{vit_small_patch14_reg4_dinov2} with
\path{pretrained=True}, \path{img_size=512}, and
\path{patch_size=16} performs both the patch-embedding kernel
resize and the position-embedding resize described in
\Cref{sec:method-backbone}:

\begin{quote}
\small
\texttt{Resize patch embedding torch.Size([384, 3, 14, 14])}\\
\texttt{to (16, 16), w/ bicubic interpolation.}\\
\texttt{Resized position embedding: (37, 37) to (32, 32).}
\end{quote}

\textbf{HRNet-W18.} Runs in a separately resolved virtual
environment, verified to have resolved to Python 3.12.3,
PyTorch 2.7.0+cu126, Torchvision 0.22.0+cu126, and CUDA 12.6.

\textbf{RTMPose-s.} Runs in a third dedicated environment: Python
3.10, PyTorch 2.1.0+cu121, Torchvision 0.16.0+cu121, MMCV 2.1.0,
MMEngine 0.10.7, MMDetection 3.2.0, and MMPose 1.3.2 (pinned commit
\texttt{5408bc76f5b8}). Additional runtime compatibility settings,
including setuptools pinning and deterministic-CUDA configuration,
are recorded in the archived environment file and each run's
provenance JSON.


% =================================================================
% A.2 Data Ethics, Licensing, and Provenance
% =================================================================

\section{Data Ethics, Licensing, and Provenance}
\label{app:ethics}

The original sources report separate ethical approval and governance
arrangements for each fetal ultrasound cohort. They are cited
separately because \textcite{divece2025multicentre} does not provide
a single ethics statement covering all three sources.

\textbf{UCL.} Data collection was approved under IRAS ID 230125
\parencite{divece2025multicentre}.

\textbf{FP.} Data collection was approved by the coordinating site's
Institutional Review Board (Comit\'e de {\'e}tica de investigaci{\'o}n
cl\'inica, ID HCB 2018/0031), and every participant provided written
informed consent \parencite{burgosartizzu2020fetalplanes}.

\textbf{HC18.} Data collection was approved by the CMO
Arnhem-Nijmegen regional ethics committee. For this retrospective
collection, consent was waived and the images were anonymised in
accordance with the Declaration of Helsinki
\parencite{vandenheuvel2018hc18}.

\textbf{Licensing.} The published journal article's CC~BY~4.0
licence covers its text and figures only, not the underlying image
data. The UCL Research Data Repository record hosting the pooled
Multicentre release is licensed under CC~BY-NC-SA~4.0, confirmed
directly against that record's own Figshare and DataCite metadata
\parencite{divece2025rdr}. The data was obtained as the consolidated
pooled release following the procedure documented by the official
code repository \parencite{multicentrefetalbiometryrepo}. The
underlying datasets are not redistributed with this thesis. The
de-identified UCL test images used as annotated qualitative examples
are reproduced under this licence; the applied resizing and
annotations are stated in the corresponding figure caption.


% =================================================================
% A.3 Dataset File and Grouping Checks
% =================================================================

\section{Dataset File and Grouping Checks}
\label{app:multicentre-audit}

For internal validation splitting of the pooled Multicentre training
data, FP filenames are grouped by the patient number in
\texttt{Patient[ID]\_Plane...}, while HC18 and UCL filenames are
grouped by their leading numeric prefix. An explicit source tag is
prepended to prevent equal numeric identifiers from different source
datasets being merged. The HC18 source tag is assigned by
cross-referencing the original HC18 annotation files because HC18 and
UCL use similar filename conventions.

Row counts in the released annotation files were also checked against
the constituent sources. They sum exactly across sources, with no
duplicate \texttt{image\_name} within any single released split and
no identical filename appearing in both members of a released
Train/Test pair.


% =================================================================
% A.4 Historical Evaluation Conventions
% =================================================================

\section{Historical Evaluation Conventions}
\label{app:historical-nme}

An earlier, simpler \emph{fixed-channel} NME convention matched
predicted and ground-truth landmarks in a fixed query-channel order,
without the unordered-pair minimisation of PI-NME. With
$d_{\text{std}} = \lVert \hat{p}_0 - p_0 \rVert_2 + \lVert
\hat{p}_1 - p_1 \rVert_2$ and $D = \lVert p_0 - p_1 \rVert_2$:

\begin{equation}
  \label{eq:nme-fetal-fixed}
  \mathrm{NME}_{\text{fixed}} = \frac{d_{\text{std}}}{2D}.
\end{equation}

This convention is retained only to document historical training-time
monitoring and diagnostic outputs. It is not used in any formally
reported result or cross-model comparison.

An earlier version of the evaluation code computed distances in
model-input space rather than original-image space. Because the
resize from non-square ultrasound frames to the fixed
$512 \times 512$ input is anisotropic, this could affect both the
NME magnitude and the selected landmark correspondence. The corrected
implementation maps all predictions back to original-image
coordinates before evaluating either correspondence. Every PI-NME
value reported in this thesis uses this corrected computation.


% =================================================================
% A.5 Baseline Reproduction
% =================================================================

\section{Baseline Reproduction}
\label{app:baseline-reproduction}


\subsection{HRNet-W18}
\label{app:hrnet-reproduction}

The reproduction uses the original authors' training code from the
official \path{surgical-vision/Multicentre-Fetal-Biometry}
repository, detached at commit \texttt{21ee7cd70b9a}. Two audited
compatibility patches are applied, neither altering the model,
training protocol, or evaluation:

\begin{itemize}

  \item \path{forward_compat.patch}: replaces the removed
    \texttt{np.math.floor} alias with \texttt{math.floor} and adds
    explicit \texttt{weights\_only=False} for trusted local
    checkpoints.

  \item \path{apply_controlled_seed_patch.py}: seeds Python,
    NumPy, PyTorch, and CUDA RNGs, seeds the DataLoader generator
    and workers, and enables deterministic cuDNN behaviour, giving
    HRNet-W18 the same five named seeds used throughout this thesis.

\end{itemize}

The upstream training program reads the official test partition for
its per-epoch validation pass and uses that metric to create
\texttt{model\_best.pth}; no separate internal validation partition
is provided. Evaluation uses the upstream \texttt{final\_state.pth}
convention. The final state is therefore not chosen as the epoch with
minimum test NME, but the test split was monitored throughout
training.


\subsection{HRNet-W18 Resolution and Decode-Size Adaptation}
\label{app:hrnet-512-fix}

At $512 \times 512$ input, HRNet-W18's native stride-4 output
produces a $128 \times 128$ heatmap rather than the $64 \times 64$
grid at the released $256 \times 256$ input. The target-generation
grid was adjusted to match this native output. A second issue was
identified in the released coordinate decoder, which hard-codes a
$[64, 64]$ output grid at three call sites. At $128 \times 128$
actual output, this silently applies an incorrect coordinate scale.

The fix is implemented by
\path{apply_dynamic_decode_size_patch.py}. It replaces all three
sites with \path{list(config.MODEL.HEATMAP_SIZE)} and aborts the
pipeline if any old literal remains. Neither change
alters the model architecture, loss, or optimisation.


\subsection{RTMPose-s}
\label{app:rtmpose-reproduction}

The full adaptation protocol is recorded in
\path{rtmpose_reproduction/PROTOCOL_LOCKED.md}. Key codec
settings are \path{in_featuremap_size=(16,16)},
\path{simcc_split_ratio=2.0} (1024 bins per axis), and SimCC
label sigma $(8.0,8.0)$. The sigma value is extrapolated from the
official settings via
$\sigma_{\text{axis}}=\sqrt{\text{axis length}/8}$, rather than being
a published $512 \times 512$ setting. This SimCC sigma controls
one-dimensional coordinate-label distributions and is therefore not
directly comparable to the Gaussian heatmap sigma used by HRNet-W18
or QCLand.

The locked MMPose, MMEngine, MMCV, and MMDetection versions and the
official CSPNeXt-s pretrained-weight checksum are recorded alongside
the protocol. Each run's actual installed versions are additionally
captured in that run's own provenance JSON.


% =================================================================
% A.6 Four-Model Statistical Reproduction
% =================================================================

\section{Four-Model Statistical Reproduction}
\label{app:paired-stats}

The frozen four-model comparison is produced by a single invocation
of \path{freeze_four_model_comparison.py}, which iterates over
ten dataset--task cells. For each cell, the script reads per-image
PI-NME values from the audited canonical per-image CSVs (QCLand and
HRNet-W18) or computes them from per-seed files (RTMPose-s).
Agreement checks using \path{math.isclose(..., abs_tol=1e-6)}
abort the pipeline on disagreement.

The common image subset for each cell is the filename intersection
across all four methods. For each method and image, PI-NME is first
averaged across the five training seeds. Pairwise differences are
then formed between these per-image five-seed means on the common
image subset. The paired bootstrap (20{,}000 replicates, seed
20260812) reports the 2.5th/97.5th percentiles of the resampled mean
differences. The Wilcoxon test uses
\texttt{scipy.stats.wilcoxon} with \texttt{zero\_method="pratt"}.
Holm correction is applied both within each cell's six comparisons
and globally across all sixty (reported as Holm-60).

The script writes three principal outputs under
\path{four_model_freeze_20260812_v1/}:
\path{four_model_main_table.tsv} (source for the main
comparison table), \path{four_model_paired_statistics.tsv}
(source for the paired contrast tables), and
\path{four_model_asset_audit.tsv} (guard outcomes and image
counts). Per-cell per-image CSVs and a SHA-256 manifest are also
archived.


% =================================================================
% A.7 Implementation Audits for QCLand
% =================================================================

\section{Implementation Audits for QCLand}
\label{app:method-implementation-detail}


\subsection{DINOv3 Interface Adaptation}
\label{app:dinov3-adapter}

The shared backbone interface is written against \texttt{timm}'s
attribute naming, which DINOv2 exposes natively. DINOv3, loaded
through HuggingFace \texttt{transformers}, uses different attribute
names. The adapter \texttt{ViT.transformers\_to\_timm} in
\texttt{models/vit.py} remaps the loaded module onto the expected
interface, aliasing \texttt{embeddings} to \texttt{patch\_embed}
and \texttt{layer} to \texttt{blocks}, and computing
\texttt{num\_prefix\_tokens} from the register-token configuration.


\subsection{Coordinate Recovery}
\label{app:coordinate-space-fix}

Producing the original-image-space coordinates used by every QCLand
PI-NME value required two coordinate-space corrections. First, the
per-image coordinate export uses a naive multiplication to convert
from heatmap to model-input space, which is not the exact inverse of
the pixel-centre-aligned encoding; this leaves a constant additive
offset ($-3.5$ pixels for the $512/64$ configuration) that is
corrected before evaluation. Second, converting from model-input
space to original-image space requires the pixel-centre-aligned
inverse of \Cref{eq:udp-scale}, not a plain proportional rescaling.
Both corrections are applied before landmark matching or scoring.


\subsection{DeconvHeadV2 Zero Initialisation}
\label{app:deconvheadv2-init-gradient}

DeconvHeadV2's final convolution is zero-initialised. Because the
weight is exactly zero at step 0, the DeconvHeadV2 loss term sends
no gradient through it to the FiLM generator or the preceding
convolution on the first backward pass. Only the final convolution
is updated by its own loss at step 0; the FiLM generator and first
convolution begin receiving a useful gradient signal once that weight
has moved away from zero. This is a within-module, first-step effect
only: the rest of the network receives gradients from the auxiliary
losses throughout, unaffected by this initialisation.


% =================================================================
% A.8 Deterministic Training and Run Provenance
% =================================================================

\section{Deterministic Training and Run Provenance}
\label{app:seed-control}


\subsection{Seed Control}

To make data ordering and worker-level augmentation deterministic,
QCLand uses an explicit \texttt{loader\_seed} argument. When set,
each DataLoader is constructed with its own \path{torch.Generator},
seeded from \path{loader_seed} plus a fixed per-split offset
(train $+0$, validation $+1$, test $+2$). Each DataLoader also uses
\path{worker_init_fn=seed_worker}, which seeds NumPy's global
RNG inside each worker process. When \texttt{loader\_seed} is left
at its default of \path{None}, this explicit seeding is skipped,
preserving backward compatibility with older experiments. All
formally reported QCLand results were generated with
\path{loader_seed} matched to the corresponding training seed.

HRNet-W18 and RTMPose-s use their own independently audited seeding
mechanisms: the controlled-seed patch for HRNet-W18 and MMEngine's
built-in \texttt{randomness} configuration for RTMPose-s.


\subsection{Run Provenance}

QCLand, HRNet-W18, and RTMPose-s each archive per-run effective
configurations alongside their checkpoints. For QCLand, ablation
drivers generate and archive an effective YAML configuration for each
run. For HRNet-W18, the canonical YAML files and patched source are
retained. For RTMPose-s, a generated Python configuration and a
provenance JSON capturing the actual installed environment are
archived per run. The archived per-run configuration, not a static
source file, is the authoritative record of what produced each
reported checkpoint.
